% ============================================================
% 3.3.5 NORMES OR REFERENCE STANDARDS
% ============================================================

\hypertarget{p3-3-5}{}

\subsection{Reference Standards}

\subsubsection{General}

The \textbf{BN-100 LUX} Infant Radiant Warmer is designed and manufactured
with reference to applicable technical and safety requirements identified in
the manufacturer documentation available for the device.

The standards and reference documents identified in the available
documentation are presented below according to their documented scope.

No additional standard or conformity claim is introduced where the
corresponding manufacturer documentation is not available in the submission
file.

\subsubsection{IEC 60601-1}

\textbf{IEC 60601-1} is identified as the applicable general standard for the
basic safety and essential performance of medical electrical equipment.

The available manufacturer documentation indicates that compliance with
the applicable requirements of IEC 60601-1 is relevant to the electrical
safety and essential performance of the device.

The corresponding test report or certificate is to be included in the
technical documentation where available.

\subsubsection{IEC 60601-1-2}

\textbf{IEC 60601-1-2} is identified as the applicable collateral standard
concerning electromagnetic compatibility of medical electrical equipment.

The standard is relevant to the electromagnetic compatibility of the
\textbf{BN-100 LUX} and to the assessment of electromagnetic emissions and
immunity applicable to the device.

The corresponding EMC test report is to be included in the technical
documentation where available.

\subsubsection{IEC 60601-1-8}

\textbf{IEC 60601-1-8} is referenced in the available manufacturer
documentation in relation to alarm systems for medical electrical equipment.

This standard is relevant to the audible and visual alarm functions
implemented in the \textbf{BN-100 LUX}, including safety-related alarm
conditions identified in the manufacturer documentation.

\subsubsection{ISO 14971}

\textbf{ISO 14971} is the reference standard for the application of risk
management to medical devices.

The manufacturer's risk management documentation identifies hazards,
hazardous situations, risks and corresponding risk-control measures
associated with the device.

The available Risk Management Report is therefore considered as part of the
risk-management documentation supporting the safety assessment of the
device.

\subsubsection{Relationship with the Device}

The standards identified above are related to different aspects of the
\textbf{BN-100 LUX}, including:

\begin{itemize}

    \item basic electrical safety and essential performance;

    \item electromagnetic compatibility;

    \item alarm systems;

    \item risk management.

\end{itemize}

The application of these standards shall be understood in conjunction with
the applicable manufacturer documentation, test reports and certificates
included in the technical documentation.

\subsubsection{Documentation Status}

The available documentation identifies the standards relevant to the device,
but the presence of a reference to a standard does not, by itself, constitute
a complete conformity assessment.

Where a standard requires objective evidence of compliance, the applicable
test report, certificate or other supporting document should be included in
the technical documentation.

No test result or conformity statement is introduced in this section unless
supported by the corresponding manufacturer documentation.

\subsubsection{Documentary Basis}

This section is based primarily on the following manufacturer documentation:

\begin{itemize}

    \item Risk Management Report (BN-100/BN-100A), Document No.
    YF-7.3-24-02 (w), Version A/0;

    \item BN-100 / BN-100 LUX User Manual.

\end{itemize}
